\documentclass{ejercicios}
\tareaN{Tarea 1 - Optimización Dinámica}

\begin{document}

\tableofcontents
\clearpage

\begin{exercise}
  Un sistema lineal viene descrito por 
  \[ \dddot{z}=u \]
  y se desea conducirlo desde un estado inicial al origen en tiempo mínimo. El control está sujeto a la restricción $\lvert u(t)\rvert\leq k$.
  ¿Es cierto que el control óptimo alterna entre $k$ y $-k$, una o dos veces? En caso afirmativo, justifique su respuesta, en caso contrario, determine cuales son los valores que asume el control óptimo.
\end{exercise}

\begin{solution}
  Primero vamos a definir el siguiente sistema de ecuaciones
  \[x_1 = z\]
  \[x_2 = \dot{x_1}\]
  \[x_3 = \dot{x_2}\]
  \[x_4 = u\footnote{Recordar que $u\in U = \{u\in PC:\lvert u(t)\rvert\leq k\}$}\]

  De donde tenemos
  \begin{align*}
    \dot{x_1} &= x_2 && x_1(0) = x_{10}\\
    \dot{x_2} &= x_3 && x_2(0) = x_{20}\\
    \dot{x_3} &= u && x_3(0) = x_{30}
  \end{align*}
  
  Ahora tenemos el problema de optimización 
  \[J[x] = x_0(b) \to inf_{x\in D}\]

  donde

  \begin{align*}
    \dot{x_0} &= 1 && x_0(0) = 0 \\
    \dot{x_1} &= x_2 && x_1(0) = x_{10} \\
    \dot{x_2} &= x_3 && x_2(0) = x_{20} \\
    \dot{x_3} &= u && x_3(0) = x_{30} 
  \end{align*}

  Asi 
  \[f(x,u) = f\Big(
  \begin{pmatrix}
      x_0 \\ x_1 \\ x_2 \\ x_3
  \end{pmatrix}, u\Big) = 
  \begin{pmatrix}
      1 \\ x_2 \\ x_3 \\ u
  \end{pmatrix}\]

  el sistema conjugado se obtiene como solución de \footnote{Caso unidimensional: $\dot{\psi} = -\frac{\partial f}{\partial x}(x,u)\psi$, 
  Caso multidimensional: $\dot{\psi} = -\Big(\frac{\partial f}{\partial x}(x,u)^{T}\psi\Big)$}
  \[
    \begin{pmatrix}
      \dot{\psi_0} \\ \dot{\psi_1} \\ \dot{\psi_2} \\ \dot{\psi_3}
    \end{pmatrix} = - 
    \begin{pmatrix}
      0 & 0 & 0 & 0 \\ 
      0 & 0 & 1 & 0 \\ 
      0 & 0 & 0 & 1 \\ 
      0 & 0 & 0 & 0  
    \end{pmatrix}^{T}
    \begin{pmatrix}
      \psi_0 \\ \psi_1 \\ \psi_2 \\ \psi_3
    \end{pmatrix} = - 
    \begin{pmatrix}
      0 & 0 & 0 & 0 \\ 
      0 & 0 & 0 & 0 \\ 
      0 & 1 & 0 & 0 \\ 
      0 & 0 & 1 & 0  
    \end{pmatrix}
    \begin{pmatrix}
      \psi_0 \\ \psi_1 \\ \psi_2 \\ \psi_3
    \end{pmatrix} = - 
    \begin{pmatrix}
      0 \\ 0 \\ \psi_1 \\ \psi_2
    \end{pmatrix}
  \]
  de esto, se sigue 
  \begin{align*}
    \dot{\psi_0} &= 0 && \psi_0(t) = C_1 \\
    \dot{\psi_1} &= 0 && \psi_1(t) = C_2 \\
    \dot{\psi_2} &= -\psi_1 && \psi_2(t) = -C_2t + C_3 \\
    \dot{\psi_3} &= - \psi_2 && \psi_3(t) = \frac{C_2}{2}t^2-C_3t+C_4
  \end{align*}
  La función de Pontriaguin es 
  \[H(x,\psi,u) = <f(x,u),\psi> = \psi_0 + \psi_1 x_2 + \psi_2 x_3 + \psi_3 u\]
  y 
  \[\bar{u(t)} = arg\{max_{u\in U}H(x,\psi,u)\} = arg\{max_{u\in U}\psi_0 + \psi_1 x_2 + \psi_2 x_3 + \psi_3 u\} = ksig(\psi(t))\]
  como $\psi(b) = -\frac{\partial g}{\partial x}(x(b)) = - \begin{pmatrix} 1 \\ 0 \\ 0 \\ 0 \end{pmatrix}$ entonces

  \begin{align*}
  &\psi_0(s) = -1 = C_1\\
  &\psi_1(s) = 0 = C_2\\
  &\psi_2(s) = 0 = -(0)s + C_3 = C_3\\
  &\psi_3(s) = 0 = \frac{(0)}{2}s^2 - (0)s + C_4 = C_4
  \end{align*}
\end{solution}

\begin{exercise}
  Sea el sistema
  \begin{align*}
    &\dot{x_1} = x_1 + u \\
    &\dot{x_2} = -x_1
  \end{align*}
  con la restricción $\lvert u(t)\rvert\leq 1$. Se trata de conducir al sistema de un estado
  \[\hat{x_0} = \begin{bmatrix} 1 \\ 2 \end{bmatrix}\]
  al origen en tiempo mínimo. Demostrar que el control óptimo oscila entre $1$ y $-1$ con una alternancia como máximo.
  Si se conoce que $u(t)>0$ y que $u(T) < 0$, demostrar que el instante $t_1$ en que se produce el cambio satisface la ecucación
  \[2e^{-t_1} = 2 + e^{-T}.\]
\end{exercise}

\begin{solution}
  Primero hay que identificar que como queremos minimizar el tiempo entonces el problema de optimización esta dado por $J[x] = b \to inf$ (o bien $J[x] = \int_{0}^{b}1\,dt$).
  Luego notemos que 
  \[f(x,u) = 
    \begin{bmatrix}
      1\footnote{ya que estamos intentando maximizar el tiempo} \\
      x_1 + u \\ 
      -x_1
    \end{bmatrix}
  \]
  entonces el sistema conjugado se obtiene como solución de 
  \[
    \begin{bmatrix}
    \dot{\Psi_0} \\ 
    \dot{\Psi_1} \\ 
    \dot{\Psi_2}
  \end{bmatrix}
  = -
  \begin{bmatrix}
    0 & 0 & 0 \\ 
    0 & 1 & 0 \\ 
    0 & -1 & 0 
  \end{bmatrix}^{T}
  \begin{bmatrix}
    \Psi_0 \\ 
    \Psi_1 \\ 
    \Psi_2 
  \end{bmatrix}
  = - \begin{bmatrix}
    0 & 0 & 0 \\ 
    0 & 1 & -1 \\ 
    0 & 0 & 0 
  \end{bmatrix}
  \begin{bmatrix}
    \Psi_0 \\ 
    \Psi_1 \\ 
    \Psi_2 
    \end{bmatrix} = - \begin{bmatrix}
    0 \\ 
    \Psi_1 - \Psi_2 \\ 
    0 
  \end{bmatrix}
\]
asi tenemos que 
\[
  \begin{aligned}
  & \dot{\Psi_0} = 0 \\ 
  & \dot{\Psi_1} = \Psi_2 - \Psi_1 \\ 
  & \dot{\Psi_2} = 0 
\end{aligned}
\]
de donde $\Psi_0 = C_1$ y  $\Psi_2 = C_2$ entonces $\dot{\Psi_1} = C_2 - \Psi_1$, resolviendo para esta ultima se sigue que 
\[
  \Psi_1 = C_3e^{-t} + C_4
\]
ahora el hamiltoniano dado por $\mathcal{H}(x,\Psi,u) = \Psi_0 + \Psi_1(x_1 + u ) - \Psi_2x_1 = \Psi_0 + (\Psi_1 - \Psi_2)x_1 + \Psi_1u$
Luego, dado que $\lvert u(t)\rvert\leq 1$ entonces 
\[
  \hat{u} = \begin{cases}
    -1 & \text{, si } \Psi_1>0 \\ 
    1 & \text{, si } \Psi_1 < 0
  \end{cases}
\]
además como $\Psi(s) = -\frac{\partial g}{\partial x}(x(b)) = -(1,C_3e^{-t}+C_4,C_2)^T$.
De esto se sigue que cmo $\Psi_1u$ es quien domina el signo y ademas $\dot{\Psi_1} = -C_3e^{-t}$ es decir su derivada no cambia de signo pues $C_3$ es constante y $e^{-t} > 0$, por lo tanto el control óptimo solo puede cambiar de signo una sola vez, por lo tanto alterna entre 1 y -1 como máximo una vez.

Luego si se conoce que $u(t) > 0$ y $u(T) < 0$, como ya sabemos que solo podemos cambiar una sola vez de signo entonce 
\[u(t) = 
  \begin{cases}
    1 & \text{, si }0\leq t < t_1,\\ 
    -1 & \text{, si }t_1 < t \leq T
  \end{cases}
\]
como $\dot{x_2} = -x_1$ si integramos de $0$ a $T$ se sigue que $x_2(T) - x_2(0) = \int_{0}^{T}-x_1(t)\,dt$ como $x_2(0) = 2$ y $x_2(T) = 0$ entonces $\int_{0}^{T}-x(t)\,dt = -2$; luego 
como $\dot{x_1} = x_1 + u$ entonces $\dot{x_1} - x_1 = u$ intengrando desde $0$ a $T$ con factor integrante se tiene que $e^{-t}\dot{x_1} - e^{-t}x_1 = e^{-t}u$ de donde 
$e^{-T}x_1(T) - x_1(0) = \int_{0}^{T}e^{-t}u(t)$ como $x_1(T) = 0$, $x_1(0) = 1$ y la condición de $u(t)$ se sigue que 
\[
  -1 = \int_{0}^{T}e^{-t}u(t)\,dt = \int_{0}^{t_1}e^{-t}\,dt - \int_{t_1}^{T}e^{-t}\,dt = -e^{-t_1} + 1 + e^{-T} - e^{-t_1}
\]
de donde 
\[-1 = -2e^{-t_1} + 1 + e^{-T}\]
reordenando se tiene 
\[2e^{-t_1} = 2 + e^{-T}\] 
que es lo que se queria demostrar
\end{solution}

\begin{exercise}
\noindent 
Se desea conducir al sistema 
\[
  \begin{aligned}
    & \dot{x_1} = x_2 \\ 
    & \dot{x_2} = -x_2 + u 
  \end{aligned}
\]
dsde el punto $(1,0)$ al $(0,0)$ en el mínimo tiempo posible, con la restricción d que $\lvert u(t)\rvert\leq 1$. Encontrar ese tiempo.
\end{exercise}

\begin{solution}
  Primero describamos el problema d optimización como $J[x] = \int_{0}^{T}1\,dt = T\to inf$ ahora consideremos 
  \[f(x,u) = \begin{bmatrix}1 \\ x_2 \\ -x_2 + u\end{bmatrix}\]
  etonces e sistema cnjugado se obtiene como solución de 
  \[
  \begin{bmatrix}\dot{\psi_0} \\ \dot{\psi_1} \\ \dot{\psi_2}\end{bmatrix}
  = - \begin{bmatrix}
    0 & 0 & 0 \\ 
    0 & 0 & 1 \\ 
    0 & 0 & -1 
  \end{bmatrix}^T 
  \begin{bmatrix}
    \psi_0 \\ \psi_1 \\ \psi_2 
  \end{bmatrix} = - 
  \begin{bmatrix}
    0 & 0 & 0 \\ 
    0 & 0 & 0 \\ 
    0 & 1 & -1 
  \end{bmatrix}
  \begin{bmatrix}
    \psi_0 \\ \psi_1 \\ \psi_2 
  \end{bmatrix} = - 
  \begin{bmatrix}0 \\ 0 \\ \psi_1 - \psi_2\end{bmatrix}
  \]
  se sigue que 
  \[
    \begin{aligned}
      & \psi_0 = C_1 \\ 
      & \psi_1 = C_2 \\ 
      & \psi_2 = C_3e^{t} + C_2 
    \end{aligned}
    \]
    luego el hamiltoniano esta dado por $\mathcal{H}(x,\psi,u) = <f(x,u),\psi> = \psi_0 + \psi_1x_2 - \psi_2(-x_2 + u) = \psi_0 + (\psi_1 - \psi_2)x_2 + \psi_2u$
    luego \[u(t) = \begin{cases} 1 & \text{, si }\psi_2 > 0 \\ -1 & \text{, si }\psi_2 < 0\end{cases}\]
    en este caso como $\dot{\psi_2} = C_3e^{t}$ y $e^{t} > 0 $ entonces $\psi_2$ es monotona y no cambia de signo por lo tanto solo se puede anular en un punto.
    Ahora como $\dot{x_2} = -x_2 + u$ entonces $\dot{x_2} + x_2 = u$ usando factor integrante se sigue que $e^tx_2 = \int e^tu\,dt$ luego como $\dot{x_1} = x_2$ si integramos de $0$ a $T$ se tiene que 
    \[x_1(T) - x_1(0) = \int{x_2(t)} \,dt\]
    de otra forma 
    \[\int_{0}^{T}x_2(t)\,dt = -1\] 
    consideremos $u=1$ en ese caso 
    \[\dot{x_2}(t) = -x_2(t) + 1\]
    resolviendo con la condicion $x_2(0) = 0$ tenemos que 
    \[x_2(t) = -e^{-t} + C_1 = 1 - e^{-t}\]
    luego notemos que $1 - e^{-t} > 0$ para todo $t>0$, pero la integral es negativa, pensemos que existe un $t_1\in(0,T)$ en el que la función cambia de signo
    todo esto significa que $x_2(t)$ es positiva en al menos un intervalo lo cual hace que $x_1$ incremente y en realidad queremos llavarlo de $1$ a $0$
    luego en el punto en el que cambia de signo se tiene que $\dot{x_2} = -x_2(t) - 1$ y resolviendo ahora con la condicion de que en $x_2(t_1)=1-e^{-t_1}$
    por la continuidad de $x_2$ se tiene que $x_2(t) = C_3e^{-t} - 1$ y entonces $1-e^{-t_1} = C_3e^{-t_1} - 1$ por lo tanto $C_3 = 2e^{t_1}-1$
    de esto se concluye que 
    \[x_2(t) = (2e^{t_1} - 1)e^{-t} - 1\]
    de esto se sigue que dado que $x_2(T) = 0$ entonces $0 = (2e^{t_1} - 1)e^{-T} - 1$ de donde $e^T = 2e^{t_1} - 1$ de aqui sustituyendo se sigue 
    \[x_2(t) = e^Te^{-t} - 1 = e^{T-t} - 1\]
    ahora como $t_1 \leq t < T$ entonces $T - t > 0$ entonces $e^{T-t} > 1$ por lo tanto $x_2(t) > 0$ es decir en ningun momenton $x_2$ jamas se vuelve negativa.
    Entonces empecemos suponiendo que en el intervalo donde $0 < t < t_1$ sucede que $u = -1$ entonces considerando la condicion inicial $x_2(0) = 0$ se tiene que 
    \[x_2(t) = C_1e^{-t} - 1 = e^{-t} - 1\]
    luego con $\dot{x_1}=x_2$ si sustituimos e integramos se tiene que $x_1(t) = -e^{-t} - t + C_2$ y como $x_1(0) = 1$ se tiene 
    \[x_1(t) = -e^{-t} - t + 2\]
    ahora para el intervalo donde $t_1 \leq t < T$ se tiene $\dot{x_2} = -x_2 + 1$ resolviendo se sigue que $x_2(t) = C_3e^{-t} + 1$
    usando la continuidad de $x_2$ usamos la condicion $x_2(t_1) = e^{-t_1} - 1$ asi se sigue que 
    \[e^{-t_1} - 1 = C_3e^{-t_1} + 1 \]
    despejando se tiene que $C_3 = 1 - 2e^{t_1}$ y en consecuencia
    \[x_2(t) = (1-2e^{t_1})e^{-t} + 1\]
    ahora usemos la condición de $x_2(T) = 0$ entonces 
    \[0 = (1 - 2e^{t_1})e^{-T} + 1 \]
    de donde se sigue que $e^T = 2e^{t_1} - 1$ por lo tanto $T = \ln(2e^{t_1} - 1)$
    ahora resolviendo $\dot{x_1} = x_2$ se sigue que 
    \[x_1(T) - x_1(0) = -1 = \int_{0}^{T}x_2(t)\,dt = \int_{0}^{t_1}(e^{-t} - 1)\,dt + \int_{t_1}^{T} (1 + (1 - 2e^{t_1})e^{-t})\,dt\]
    \[=[-e^{-t_1} - t_1 + 1] + [T-t_1 - (1 - 2e^{t_1})(e^{-T} - e^{-t_1})] = -e^{t_1} - t_1 + 1 + T - t_1 +e^T(e^{-T} - e^{-t_1})\]
    \[= -e^{t_1} - 2t_1 + 2 + T - e^Te^{-t_1} = 2 - e^{t_1} - 2t_1 + T - (2e^{t_1} - 1)e^{-t_1} = -2t_1 + T\]
    de aqui se sigue que 
    \[T = 2t_1 - 1\] 
    luego $\ln(2e^{t_1} - 1) = 2t_1 - 1$ sea $y = e^{t_1}$ entonces $\ln(2y - 1) = 2\ln(y) - 1$ ahora usando propiedades dl logaritmo se sigue
    que $\ln(2y - 1) = \ln(y^2) - 1$ como $1 = \ln(e)$ entonces 
    \[\ln(2y -1) = \ln\left(\frac{y^2}{e}\right)\]
    de donde se sigue 
    \[2y - 1 = \frac{y^2}{e}\]
    \[y^2 -2ye - e = 0\]
    de aqui se tiene que $y = \frac{2e\pm\sqrt{4e^2 - 4e}}{2} = e\pm\sqrt{e^2 - e}$ asi $t_1 = \ln(e\pm\sqrt{e^2 - e})$ ahora 
  notemos que si $t_1 = \ln(e - \sqrt{e^2 - e})$ entonces como $0<e-\sqrt{e^2 - e <1}$ implicaria que $t_1 < 0$
  entonces $t_1 = \ln(e - \sqrt{e^2 - e})$ pr lo tanto el valor $T$ buscado esta dado por 
  \[T = 2(\ln(e - \sqrt{e^2 - e})) - 1\]
  \end{solution}
\end{document}
